\section{Introduction}
\label{sec:introduction}

\subsection{The Seed-Invariance Problem}

The B-series of algorithmic redistricting papers has converged on a single
legislative vehicle: the \dia{}, which proposes to mandate the \ar{}
algorithm---a recursive bisection procedure derived from the \hh{} priority
rule---as the exclusive method for congressional redistricting in all fifty
states~\citep{b02paper}.
The statute's reproducibility guarantee is a seed: anyone who runs the
reference implementation with the published seed and the public TIGER/Line
inputs must obtain the same map.

That design has a gap.

METIS---the graph-partitioning engine underlying \ar{}---is a heuristic.
Graph minimum bisection is $\mathcal{NP}$-hard~\citep{b6paper}; METIS finds a
local optimum via multilevel Kernighan-Lin refinement, not the global minimum.
The seed controls which local optimum is reached.
A fixed seed is reproducible; it is not guaranteed to be optimal.

The practical consequence was documented in B.7~\citep{b7paper}.
Running all fifty states over a sweep of 1{,}500 seeds each, the analysis
found that for most states the algorithm converges quickly---subsequent seeds
produce no improvement after a short run---but for a handful of states,
including Georgia and Wisconsin, the last improvement occurs at surprisingly
high seed offsets (seed 489 and seed 1{,}023 relative to the content-derived
start, respectively).

A statute that publishes a single fixed seed cannot certify that the map it
produces is the best available map under its own objective function.
A well-prepared challenger need only run a different seed and show a
lower edge cut to challenge the map's optimality claim.

\subsection{The ConvergenceSweep Solution}

\cs{} eliminates the single-seed fragility.
Instead of fixing one seed, the algorithm sweeps forward from the
content-derived seed $s_0$ and halts when $T$ consecutive seeds produce no
improvement in the normalised edge cut $\ECnorm$.
The final map is the one with the lowest $\ECnorm$ seen across all seeds tried.

The design achieves two properties simultaneously:
\begin{enumerate}
\item \textbf{Determinism.} Given the census release identifier and the
      protocol constant \texttt{"DIA\_SEED\_V1"}, $s_0$ is uniquely determined.
      Given $T$ and $s_0$, the sweep produces the same map on every machine.
      The result is fully auditable.
\item \textbf{Optimality.} If the algorithm has not improved in $T$ consecutive
      seeds, the probability of finding a strictly better plan at seed
      $s_0 + T + k$ decreases exponentially in $k$ (Theorem~\ref{thm:gumbel}).
      For $T = 600$, the probability of missing a better plan is below
      $10^{-3}$ for every state in the union.
\end{enumerate}

Neither property holds for the fixed-seed design.

\subsection{Paper Structure and Main Results}

Section~\ref{sec:algorithm} gives the formal pseudocode for \cs{},
defines $T$-stability, and analyzes computational complexity.
Section~\ref{sec:sufficiency} presents the core empirical argument: a
50-state table of convergence tails, identification of Georgia as the
worst-case state (tail = 511), and the Gumbel extreme-value argument showing
$P(\text{tail} > 600) < 0.001$.
Section~\ref{sec:sha256} explains why the specific seed formula
$\seedfmt{\texttt{census\_release\_id} \,\|\, \texttt{"DIA\_SEED\_V1"}}$
closes the political-manipulation attack surface.
Section~\ref{sec:implementation} documents the existing \texttt{redist}
binary implementation and the federal-statute invocation string.
Section~\ref{sec:conclusion} closes by situating B.16 in the B.02
constitutional argument.

\subsection{Relation to Prior Work}

B.7~\citep{b7paper} is the empirical source for the state-level convergence
data used throughout.
B.11~\citep{b11paper} established zero seed variance for the \ar{}
prime-factorisation tree; \cs{} addresses the residual variance in the
GeoSection-based seed sweep that B.02 relies on.
B.02~\citep{b02paper} is the statutory vehicle; this paper supplies the
mathematical foundation for the seed-derivation clause in its operative text.
B.6~\citep{b6paper} established the $\mathcal{NP}$-hardness context that
makes convergence certification necessary rather than merely convenient.
